\ficha{Flandreau e Flores (2012), Bondholders versus bond-sellers}{flandreauflores2012}

\begin{identificacao}
FLANDREAU, Marc; FLORES, Juan H. Bondholders versus bond-sellers? Investment
banks and conditionality lending in the London market for foreign government
debt, 1815-1913. \emph{European Review of Economic History}, v. 16, n. 4,
p. 356-383, 2012. DOI: 10.1093/ereh/hes005.

Artigo de periódico, 28 páginas, lido integralmente, em inglês. Os apêndices A
e B, com o modelo formal e o cálculo do desempenho das apólices do funding,
existem apenas como material suplementar online e não foram consultados.
\end{identificacao}

\bloco{Problema e tese}
Quem obrigava um governo soberano a pagar, no mercado de Londres entre 1815 e
1913? A literatura macroeconômica recente responde que foram os comitês de
portadores de títulos, em especial a Corporation of Foreign Bondholders (CFB),
criada em 1868. Flandreau e Flores invertem a resposta. Quem disciplinava o
devedor era o banco emissor de prestígio, dono de um sinal de reputação que
funcionava como colateral do papel que vendia. Monitorar o cliente gerava renda,
e por isso o banco tinha ao mesmo tempo o meio e o interesse de exigir ajuste
macroeconômico como contrapartida de novo financiamento. O empréstimo
condicionado é, nessa leitura, investimento na marca do próprio banco. Os
comitês de portadores tinham punição, mas não tinham crédito a oferecer, e o
mercado não reagiu à criação da CFB.

\chave{Conditionality lending, we argue, is thus an investment in the bank's own brand.}{p. 367}

\bloco{Argumento}
\begin{enumerate}[leftmargin=*,nosep]
\item A relação não é bilateral, e sim triangular: portadores, governos e bancos
  que colocam o papel no mercado primário. Investidores comuns não distinguem
  bom de mau devedor, intermediários informados distinguem, e o mercado se
  organiza em pirâmide, com casas de prestígio monopolizando o papel de primeira
  linha e emissores comuns disputando o papel de risco.
\item Teste de confiabilidade. Os autores tomam o universo das emissões de 1850
  a 1873 e seguem o histórico de pagamento de cada título até 1878, comparando o
  retorno prometido na emissão com o retorno efetivamente realizado. A
  correlação entre os dois é superior a 0,7 nas emissões dos Rothschild e
  negativa, cerca de menos 0,5, nas demais (p. 363). O prestígio reduzia
  incerteza, e não apenas risco.
\item Banco de relacionamento. Países que trocavam pouco de banqueiro pagavam
  spreads menores, em dados de emissões entre 1877 e 1913 (p. 366). E a fatia de
  mercado dos Rothschild encolhia nos períodos de expansão e se recuperava
  depois do colapso, entre 1820 e 1900 (p. 367). O prestígio é contracíclico:
  vale pouco quando ninguém tem medo.
\item Mecanismo de imposição. A exigência de ajuste é crível porque a receita da
  casa de prestígio vem da fatia de mercado em empréstimos bem-sucedidos, não das
  comissões de cada emissão isolada. Não há dilema entre comissão e reputação, e
  quanto maior o zelo pela marca, mais duro o ajuste exigido (p. 368). Recusado o
  ajuste, o país cai para um banqueiro de segunda linha, com termos piores, e a
  própria troca sinaliza problema ao mercado.
\item O funding loan brasileiro de 1898 é o caso escolhido para observar o
  mecanismo fora do equilíbrio. A expansão monetária alimentou a depreciação
  cambial, que virou crise fiscal pelo descasamento entre dívida em libras e
  receita em moeda depreciada. Em 1898, pagar o serviço em libras teria
  absorvido 62 por cento da receita total do governo, e em maio o mercado já
  esperava moratória (p. 370). O episódio é especial porque o Brasil era país dos
  Rothschild: a tabela 1 mostra que a casa subscreveu a totalidade das emissões
  do governo brasileiro entre 1852 e 1913, e o spread de emissão cai de 3,7 em
  1893 e 3,5 em 1895 para 2,4 em abril de 1906 e 1,4 em maio de 1910 (p. 370).
\item Barings cumpre função distinta. Enquanto os Rothschild evitavam o calote
  nas próprias emissões, Barings entrava no mercado de papel podre alheio, como
  agência de cobrança que restaurava crédito de terceiros para lustrar o próprio
  escudo. O caso trabalhado é a crise dos estados norte-americanos dos anos 1840
  (p. 373 a 375).
\item Teste final. Os autores olham a reação de preço, e não a taxa de
  recuperação. A criação da CFB, em novembro de 1868, não moveu o preço dos
  títulos de Egito, Colômbia e Venezuela (p. 376). Quando a CFB agiu sozinha, o
  mercado não reagiu; quando pôde anunciar o apoio de Barings, o preço subiu
  (p. 379).
\end{enumerate}

O caso brasileiro merece detalhamento, porque é nele que a produção pública da
credibilidade aparece. O acordo foi anunciado em 20 de junho de 1898, quando os
Rothschild pediram ao \emph{The Times} que publicasse a correspondência do
presidente eleito Campos Salles, com o compromisso de estabilização fiscal
(p. 371). Os termos vieram a público em seguida, entre outros no \emph{The
Economist} de 30 de junho de 1898 (p. 371). O esquema tinha lado macroeconômico
e lado financeiro: as apólices do funding pagariam o serviço da dívida, e sua
emissão ficava condicionada ao aperto monetário, com o Brasil obrigado a
incinerar papel-moeda \emph{pari passu} com os recursos recebidos, sob
supervisão de bancos fiduciários (p. 371). Os autores registram que a
historiografia brasileira é uniformemente crítica do efeito deslocador desse
programa, e leem essa crítica como confirmação de sua própria tese, já que o
alvo primeiro do esquema era a reputação dos Rothschild (p. 371). O custo do
arranjo é medido de duas formas: a valorização imediata da carteira do funding,
entre junho e setembro de 1898, chega a 37 por cento, contra poucos pontos
percentuais nas ofertas iniciais típicas (p. 371); e os títulos de 5 por cento
de 1895 renderam cerca de 8 por cento ao ano contra 5,9 por cento prometidos, um
excedente de 210 pontos-base pago aos portadores como prêmio por terem confiado
na casa (p. 372). Daí a conclusão contraintuitiva do artigo, contra Weller
(2009): os Rothschild não socorreram o Brasil, puniram o Brasil, e o Brasil
aceitou porque a alternativa era o banimento do acesso certificado ao mercado
(p. 372).

\chave{The specific pricing of the Funding Loan underscores Rothschilds' main concern: not Brazil's welfare but Rothschilds' ability to certify other borrowers in the future.}{p. 372}

\bloco{Conceitos e definições operacionais}
\emph{Conditionality lending}: financiamento cuja liberação é condicionada a
ações de política econômica exigidas pelo banco emissor, sustentada pela ameaça
crível de retirar a certificação, e cujo alvo real é a reputação do banco, não o
bem-estar do país. \emph{Prestígio}: propriedade de um sinal de reputação que
permite certificar um devedor e sustentar fatia de mercado; opera como colateral
e produz monopólio natural, porque quem tem muito a perder não mente.
\emph{Monitoramento delegado}, pelo banco, contra \emph{monitoramento direto},
pela assembleia de portadores: o primeiro é superior quando há assimetria de
informação e problema de ação coletiva. \emph{Turnover}: probabilidade de que
dois empréstimos consecutivos de um mesmo país sejam emitidos por líderes
inteiramente diferentes, medida como número de trocas dividido pelo número de
emissões (p. 366). \emph{Linha de confiabilidade}: diagonal de 45 graus em que
retorno esperado iguala retorno realizado; abaixo dela fica a ``área de
decepção'' (p. 363). \emph{Hold-up}: captura do devedor pelo monitor delegado,
favorecida pela assimetria de informação que cresce nas crises (p. 372).

\bloco{Evidência e método}
Base própria de emissões soberanas em Londres, com preço de emissão,
características do empréstimo e histórico de pagamento. Retorno esperado
estimado como yield ao vencimento na emissão, retorno realizado como taxa
interna de retorno. Três recortes: 1850 a 1873 para a confiabilidade das
emissões, 1877 a 1913 para a relação entre turnover e spread, 1820 a 1900 para a
fatia de mercado em booms e crises. Estudo de evento em torno de datas de
anúncio, para medir poder de sinalização da CFB e de Barings. No caso
brasileiro, imprensa financeira londrina (\emph{Investors Monthly Manual},
\emph{The Times}, \emph{The Economist}) e o Rothschild Archive, de onde vem a
citação do \emph{United States of Brazil Funding Scheme} (p. 371, nota 13). No
caso norte-americano, um indicador qualitativo de esforço construído a partir da
linguagem de McGrane (1935) e Hidy (1949), somado ao ``London window'' das
edições de 1838 e 1856 do \emph{Fortune's Epitome}.

\bloco{Agentes e vertentes}
N. M. Rothschild e Sons como emissor de papel de primeira linha e único
subscritor da dívida externa federal brasileira no período coberto pela tabela
1. Baring Brothers como número dois e restaurador de crédito alheio.
Bischoffsheim e Goldschmidt como o vilão do ciclo de 1870, cujo chefe mandou
atestado médico em vez de comparecer ao Select Committee de 1875. Nathaniel de
Rothschild e Sir Henry James no depoimento de 1875. Campos Salles como o devedor
que aceita a condicionalidade. Bancos fiduciários do funding: London and River
Plate Bank, London and Brazilian Bank e Brazilianische Bank für Deutschland
(p. 371, nota 13). No campo historiográfico, os autores se opõem a Eichengreen e
Portes, Wright, Mauro e Esteves, que atribuem eficácia aos comitês de
portadores, e a Weller (2009), para quem o conflito de interesses teria levado
os Rothschild a ser complacentes com o Brasil. Alinham-se aos historiadores de
empresa Hidy, Gille e Suzuki. Citam, sem endossar, a historiografia brasileira
crítica do programa de estabilização, Barroso (1936), Fritsch (1988) e Topik
(1987).

\bloco{Críticas e limites}
O próprio texto reconhece que médias de duração de moratória são frágeis e que
resultados desse tipo dificilmente seriam falseáveis (p. 376). O estudo de
evento sobre a CFB usa três países, e os autores admitem que o efeito da criação
podia estar antecipado desde 1866 (p. 376, nota 15). O indicador de esforço de
Barings é codificação da linguagem de duas obras secundárias, não medida
independente. O artigo abstrai de propósito do controle imperial e das
instituições domésticas do devedor (p. 357), de modo que nada diz sobre a
política interna que produziu o ajuste, que é justamente o terreno da
dissertação. A tese de que houve punição, e não salvamento, é sustentada por
retorno de portador, não por efeito sobre o país; os autores admitem que a
política imposta não precisa ser boa para o crescimento. Weller oferece a
leitura rival e deve ser lido em paralelo. Sobre a Caixa de Conversão o texto
não diz nada: a tabela 1 chega a 1913, mas as emissões posteriores a 1898
aparecem apenas como preço, sem discussão do regime cambial de 1906.

\begin{dialogo}
\item Procurar nos jornais, de 1906 em diante, a invocação do funding de 1898 e
  da casa Rothschild como limite ao debate: ``compromissos'', ``credores'',
  ``crédito externo'', ``o funding''. Testa se a condicionalidade imposta em 1898
  sobreviveu como argumento doméstico contra a emissão e a favor de taxa alta,
  isto é, se o enforcement do banqueiro tinha vida política interna.
\item Procurar a citação da imprensa londrina como autoridade: \emph{The Times},
  \emph{The Economist}, telegramas de Londres, cotação do mil-réis na praça
  inglesa. O artigo mostra a credibilidade sendo produzida em público, quando os
  Rothschild pedem ao \emph{The Times} que publique a correspondência de Campos
  Salles (p. 371). A hipótese a testar é se os jornais brasileiros de 1906 a 1914
  tratavam a opinião londrina como veredito sobre a Caixa, e quais jornais mais o
  faziam.
\item Procurar o vocabulário de custódia e fiscalização na discussão da agência
  de Londres e da custódia do ouro em 1911: banco depositário, fiscalização,
  garantia. Testa se o arranjo de 1898, com bancos fiduciários supervisionando a
  incineração de papel-moeda, funcionou como modelo ou como espantalho no debate
  sobre quem guardaria o lastro da Caixa.
\item Procurar como o funding é narrado nos anos da Caixa: salvação, sacrifício
  ou negócio caro. O artigo calcula 37 por cento de valorização imediata para os
  portadores e retorno de cerca de 8 por cento contra 5,9 por cento prometidos
  (p. 371 e 372). Testa se algum jornal percebeu o preço pago, o que separaria
  crítica informada de retórica antibanqueiro.
\item Procurar indício de material pago ou inspirado por interesse bancário nas
  seções livres. O artigo documenta campanhas de imprensa e artigos pagos por
  Barings nos estados norte-americanos dos anos 1840 (p. 374), e não afirma nada
  semelhante para o Brasil. A ponte é hipótese a verificar na fonte, não achado
  transferível.
\end{dialogo}

\usonocapitulo{Parte II, como armação teórica do enforcement. Serve à afirmação
de que a disciplina externa sobre o Brasil não vinha de canhoneiras nem de
comitês de portadores, e sim do banco emissor que detinha o monopólio da
certificação, e de que o funding de 1898 foi condicionalidade, com contrapartida
monetária explícita, a incineração de papel-moeda sob supervisão de bancos
fiduciários (p. 371). Fornece o dado quantitativo direto sobre a inserção
externa brasileira no período da Caixa, o spread de emissão da tabela 1, e
fundamenta a leitura de que a credibilidade era produzida em público, na
imprensa de Londres, o que liga a parte II ao objeto empírico da dissertação.
Não sustenta afirmação alguma sobre a Caixa de Conversão em si.

\chave{The specific pricing of the Funding Loan underscores Rothschilds' main concern: not Brazil's welfare but Rothschilds' ability to certify other borrowers in the future.}{p. 372}}
